\documentclass[UTF8,a4paper,12pt]{ctexart}

\usepackage[a4paper,margin=2.5cm]{geometry}
\usepackage{graphicx}
\usepackage{float}
\usepackage{xcolor}
\usepackage{listings}
\usepackage{hyperref}
\usepackage{fancyhdr}
\usepackage{enumitem}

\hypersetup{
    hidelinks,
    unicode=true
}

\pagestyle{fancy}
\fancyhf{}
\fancyfoot[C]{\thepage}
\setlength{\headheight}{15pt}

\lstset{
    basicstyle=\ttfamily\small,
    frame=single,
    breaklines=true,
    columns=fullflexible,
    keepspaces=true,
    showstringspaces=false,
    tabsize=4,
    backgroundcolor=\color{gray!8}
}

\setlist[itemize]{leftmargin=2em}

\newcommand{\screenshot}[2]{
\begin{figure}[H]
    \centering
    \includegraphics[width=0.94\textwidth]{report_assets/practice/#1}
    \caption{#2}
\end{figure}
}

\begin{document}

\begin{titlepage}
    \thispagestyle{empty}

    \vspace*{2.6cm}

    \begin{center}
        {\LARGE\bfseries 系统开发工具基础第三周实验报告\par}
    \end{center}

    \vspace{4.2cm}

    \begin{center}
        \begin{tabular}{rl}
            姓名： & 罗越 \\
            学号： & 24160001055 \\
            日期： & 2026年9月7日
        \end{tabular}
    \end{center}

    \vspace{1.5cm}

    \begin{center}
        仓库地址：\\
        \url{https://github.com/Llj159-lab/sys-dev-tools-lab}
    \end{center}
\end{titlepage}

\pagenumbering{arabic}

\tableofcontents

\newpage

\section{实验环境与目录结构}

本周实验在 WSL 的 Bash 环境中完成，系统环境为 Ubuntu 24.04 LTS。仓库目录为：

\begin{lstlisting}[language=bash]
~/sys-dev-tools-lab/week03
\end{lstlisting}

本周实验目录结构如下：

\begin{lstlisting}[language=bash]
week03/
|-- q09/
|-- q10/
|-- q11/
|-- report_assets/
|   `-- practice/
`-- report.tex
\end{lstlisting}

本周实验主要涉及 Python 项目打包、wheel 文件生成、干净虚拟环境安装、命令行程序测试、人工检查代码差异，以及协作材料的规范化整理。

\section{第9题：从源码构建并在干净环境安装 Wheel}

\subsection*{题目要求}

在 q09 中按照给定内容创建一个最小 Python 命令行包，并验证 wheel 安装。具体要求如下：

\begin{itemize}
    \item 创建 \verb|src/greetlab/__init__.py|、\verb|cli.py| 和 \verb|pyproject.toml|，并将“学号”替换为自己的学号；
    \item 在已经安装 build 工具的课程环境中运行 \verb|python -m build|，生成 wheel；
    \item 新建一个干净虚拟环境，只从生成的 wheel 安装程序；
    \item 切换到 q09 之外运行 \verb|sdt-greet --name <学号>|，确认输出正确。
\end{itemize}

\subsection*{实验过程}

进入 q09 目录并创建源代码目录：

\begin{lstlisting}[language=bash]
cd ~/sys-dev-tools-lab/week03/q09
mkdir -p src/greetlab
\end{lstlisting}

创建 \verb|src/greetlab/__init__.py|：

\begin{lstlisting}[language=bash]
touch src/greetlab/__init__.py
\end{lstlisting}

创建 \verb|src/greetlab/cli.py|：

\begin{lstlisting}[language=Python]
import argparse

def main():
    p = argparse.ArgumentParser()
    p.add_argument("--name", required=True)
    a = p.parse_args()
    print(f"Hello, {a.name}!")
\end{lstlisting}

创建 \verb|pyproject.toml|：

\begin{lstlisting}
[build-system]
requires = ["setuptools>=68"]
build-backend = "setuptools.build_meta"

[project]
name = "greetlab-24160001055"
version = "0.1.0"

[project.scripts]
sdt-greet = "greetlab.cli:main"
\end{lstlisting}

检查目录结构：

\begin{lstlisting}[language=bash]
find . -maxdepth 3 -type f | sort
\end{lstlisting}

随后运行构建命令：

\begin{lstlisting}[language=bash]
python -m build
\end{lstlisting}

构建成功后，在 q09 目录中生成 \verb|dist| 目录，其中包含 wheel 文件。

为了验证 wheel 能够独立安装，新建干净虚拟环境：

\begin{lstlisting}[language=bash]
python -m venv .venv-clean
source .venv-clean/bin/activate
pip install dist/*.whl
\end{lstlisting}

安装完成后运行命令行程序：

\begin{lstlisting}[language=bash]
sdt-greet --name 24160001055
\end{lstlisting}

\subsection*{实验结果}

实验成功按照 \verb|src| 布局创建了 \verb|greetlab| Python 包，并使用 \verb|pyproject.toml| 配置了项目元数据和命令行入口。执行 \verb|python -m build| 后生成了 wheel 文件。

在新建的干净虚拟环境中，程序只通过生成的 wheel 安装，没有依赖原始源码目录。安装完成后执行：

\begin{lstlisting}
sdt-greet --name 24160001055
\end{lstlisting}

程序能够输出：

\begin{lstlisting}
Hello, 24160001055!
\end{lstlisting}

这表明 wheel 构建、安装和命令行入口配置均能够正常工作。

\screenshot{practice09_tree.png}{q09 项目目录结构}

\screenshot{practice09_build.png}{使用 python -m build 生成 wheel}

\screenshot{practice09_install_run.png}{在干净虚拟环境中安装并运行命令行程序}

\subsection*{分析}

本题通过 \verb|pyproject.toml| 描述 Python 项目的构建方式、项目名称、版本和命令行入口。使用 wheel 进行安装，可以验证项目是否具备独立分发能力。将源代码放在 \verb|src| 目录中，有助于避免直接从项目根目录导入未安装代码的问题。通过在干净虚拟环境中安装 wheel，可以更准确地验证打包结果，而不是依赖开发目录中的文件。

\section{第10题：让编程智能体进入可验证的修复循环}

\subsection*{题目要求}

复制 q09 为 q10，增加一个失败测试，要求命令行参数异常时程序以 \verb|SystemExit(2)| 结束。先运行测试确认失败，再向编程智能体说明目标、约束和测试命令，要求其修改实现。修改后人工检查 \verb|diff|，撤销无关修改，并再次运行测试。最后在 \verb|ai_log.md| 中记录过程，且日志不超过 5 行。

\subsection*{实验过程}

首先复制 q09：

\begin{lstlisting}[language=bash]
cd ~/sys-dev-tools-lab/week03
cp -r q09 q10
cd q10
\end{lstlisting}

创建测试目录和测试文件：

\begin{lstlisting}[language=bash]
mkdir -p tests
touch tests/test_cli.py
\end{lstlisting}

实际使用的测试代码如下：

\begin{lstlisting}[language=Python]
import pytest
from greetlab import cli

def test_main_exits_when_name_is_whitespace():
    with pytest.raises(SystemExit) as excinfo:
        cli.main()
    assert excinfo.value.code == 2
\end{lstlisting}

运行测试：

\begin{lstlisting}[language=bash]
pytest tests/test_cli.py -v
\end{lstlisting}

第一次运行测试时，原有程序因为缺少必需的 \verb|--name| 参数而退出，但测试需要捕获并验证退出状态，因此测试结果未通过。

随后向编程智能体提出如下要求：

\begin{lstlisting}
请修改 src/greetlab/cli.py，使当 --name 参数缺失时，
main() 以 SystemExit(2) 退出。

约束：
1. 只修改 src/greetlab/cli.py；
2. 不修改测试文件；
3. 不添加无关功能；
4. 修复后运行 pytest tests/test_cli.py -v。
\end{lstlisting}

编程智能体给出的修复思路是在 \verb|cli.py| 中捕获参数解析产生的 \verb|SystemExit|，并确保程序以状态码 2 结束：

\begin{lstlisting}[language=Python]
import argparse
import sys

def main():
    p = argparse.ArgumentParser()
    p.add_argument("--name", required=True)

    try:
        a = p.parse_args()
        print(f"Hello, {a.name}!")
    except SystemExit:
        sys.exit(2)
\end{lstlisting}

修改完成后再次运行：

\begin{lstlisting}[language=bash]
pytest tests/test_cli.py -v
\end{lstlisting}

然后人工查看差异：

\begin{lstlisting}[language=bash]
git diff src/greetlab/cli.py
\end{lstlisting}

最后创建 \verb|ai_log.md|：

\begin{lstlisting}[language=bash]
cat > ai_log.md <<'EOF'
# AI 修复日志
- 核心提示：要求修改 cli.py，使 --name 缺失时以 SystemExit(2) 退出
- AI 改动：在 main() 中添加参数解析异常处理
- 人工验证：运行 pytest 测试通过，diff 确认仅修改 cli.py
EOF
\end{lstlisting}

\subsection*{实际结果说明}

本题标题描述的是空白名字场景，但本次实际测试代码没有通过命令行参数设置：

\begin{lstlisting}[language=Python]
--name "   "
\end{lstlisting}

而是直接调用：

\begin{lstlisting}[language=Python]
cli.main()
\end{lstlisting}

因此，本次测试实际验证的是：当 \verb|--name| 参数缺失时，程序是否以 \verb|SystemExit(2)| 结束。报告中的测试结果和分析按照实际测试代码及截图记录。

\subsection*{实验结果}

第一次执行 pytest 时，测试未通过，证明原有程序不能满足新的退出状态要求。编程智能体根据给定的目标和约束修改了 \verb|src/greetlab/cli.py|。

人工检查 \verb|git diff| 后，确认修改集中在命令行程序实现中，没有引入与任务无关的功能。再次运行测试后，测试通过，说明当前实现能够满足本次实际测试对退出状态码的要求。

\verb|ai_log.md| 记录了提示目标、程序修改和人工验证过程，且内容控制在 5 行以内。

\screenshot{practice10_test_fail.png}{第10题第一次测试失败}

\screenshot{practice10_test_pass.png}{第10题修复后测试通过}

\screenshot{practice10_diff.png}{人工检查程序修改差异}

\screenshot{practice10_ai_log.png}{AI 修复日志内容}

\subsection*{分析}

本题体现了可验证修复循环的基本过程。首先通过测试暴露当前实现的问题，然后将目标、约束和测试命令明确传达给编程智能体，使修改范围和验收标准保持清晰。修改完成后，人工检查 \verb|git diff|，可以确认智能体没有修改无关文件或添加无关功能。最后重新运行测试，形成“失败测试、目标说明、自动修改、人工审查、再次验证”的闭环。

需要注意的是，测试名称虽然包含 whitespace，但实际代码没有构造空白字符串参数。因此本次实验结果只能证明缺少 \verb|--name| 时返回状态码 2，不能证明仅含空白字符的 \verb|--name| 已被拒绝。

\section{第11题：把“无法处理”的协作材料改成可执行信息}

\subsection*{题目要求}

围绕“空姓名仍输出问候语”这一缺陷，重写三段协作材料：

\begin{itemize}
    \item 在 \verb|communication.md| 中重写 Issue，包含环境、复现命令、期望结果和实际结果；未知信息写为“待确认”；
    \item 重写提交信息，标题使用祈使语气，正文说明问题和解决方案；
    \item 重写评审意见，指出具体行为、风险和建议动作，并标注 Blocking、Suggestion 或 Nit；
    \item 全文控制在 400 字以内。
\end{itemize}

\subsection*{实验过程}

进入 q11 目录：

\begin{lstlisting}[language=bash]
cd ~/sys-dev-tools-lab/week03/q11
\end{lstlisting}

创建或编辑 \verb|communication.md|：

\begin{lstlisting}[language=bash]
nano communication.md
\end{lstlisting}

实际整理后的协作材料如下：

\begin{lstlisting}
Issue: Windows 环境下运行 sdt-greet 时，传入仅含空白字符的 --name 仍被当作有效输入。
环境：Windows，具体版本待确认；命令版本待确认。
复现命令：sdt-greet --name "   "
期望结果：程序拒绝空白姓名，并以非 0 状态退出。
实际结果：程序输出 "Hello, !"，并以 0 状态退出。
修复建议：对 name 执行 strip，结果为空时返回参数错误。

提交信息：
fix: reject whitespace-only name

空白姓名会产生无效问候语。增加参数校验，
拒绝只包含空白字符的 name，并保留正常姓名的输出行为。

评审意见：
Blocking：当前实现接受空白姓名并输出无效结果，
建议增加校验和自动化测试。
Suggestion：补充 Windows 环境和命令版本信息。
Nit：提交说明可进一步写明测试命令。
\end{lstlisting}

检查文件长度：

\begin{lstlisting}[language=bash]
wc -m communication.md
cat communication.md
\end{lstlisting}

\subsection*{实验结果}

重写后的 \verb|communication.md| 明确记录了问题环境、复现命令、期望结果和实际结果，并将未知信息标记为“待确认”。提交信息使用了祈使式标题，正文说明了缺陷原因和修复方向。评审意见按照问题严重程度区分为 Blocking、Suggestion 和 Nit，分别对应必须修复的问题、建议补充的信息和非阻塞性的文字改进。

\screenshot{q11_communication.png}{重写后的协作材料}

\subsection*{分析}

本题将模糊的协作材料转换为可以执行和验证的问题描述。一个有效的 Issue 应当使其他开发者能够根据环境和复现命令重新观察问题，并依据期望结果和实际结果判断缺陷是否存在。提交信息需要简洁表达修改目的，评审意见则应指出具体行为、潜在风险和建议动作。通过这种结构化表达，可以减少沟通歧义，提高问题处理和代码审查的效率。

\section{解题感悟}

本周课上实验主要围绕 Python 程序的分发、测试和协作展开。第9题通过构建 wheel 并在干净虚拟环境中安装，理解了源代码、构建产物和可执行命令之间的关系。第10题通过失败测试、智能体修改、人工审查和再次测试，掌握了以测试结果约束自动化修改的方法。第11题则进一步认识到，清晰的环境信息、复现步骤、预期结果和实际结果是有效协作的基础。

通过本周实验可以看出，软件开发不仅包括编写程序，还包括打包、验证、审查和沟通。将任务目标、修改约束和验证命令明确表达，可以提高自动化工具输出的可靠性；使用干净环境、测试结果和差异检查进行复核，可以降低因环境依赖或无关修改造成的风险。

\end{document}